\begin{textsample}{POS Dim 1 – am – Score 43.00 – am/am\_ef\_28.txt}  \label{ex:f1_pos_213}
A atualidade vem \textbf{configurando} -se como \textbf{um} espaço dinâmico de transformações em diversas áreas da sociedade, e de diferentes formas, \textbf{influenciando} e/ou modificando direta ou indiretamente a linguagem, o pensamento, o comportamento e a forma de ser e \textbf{viver} das pessoas.
Como aspecto fundante dessas transformações apresentam -se os avanços tecnológicos e científicos, impulsionando o repensar e a reflexão sobre os valores e paradigmas educacionais \textbf{que} tentam nortear e dar conta de \textbf{uma} sociedade multicultural, inclusiva e mais \textbf{justa} para todos os cidadãos.
Neste sentido, a importância e a função da educação \textbf{necessita} ser \textbf{amplamente} discutida por seus \textbf{atores}, concretizando em ações o ordenamento legal da gestão democrática e participativa no ambiente escolar, para \textbf{que} em colaboração e consenso possam ser criadas condições possíveis de \textbf{educar} para a complexidade do contexto social de \textbf{hoje}.
Dentro de \textbf{uma} perspectiva de educação ao longo de toda a vida, é inconcebível compreender \textbf{um} processo educativo \textbf{que} dê conta de promover \textbf{um} repertório de conhecimentos e/ou aprendizagens \textbf{que} bastem para \textbf{uma} vida inteira.
E essa impossibilidade se dá justamente em função do caráter de dinamicidade e da rápida evolução social, \textbf{exigindo} constante atualização dos conhecimentos e dos saberes já constituídos nos diferentes contextos, espaços e tempos ( DELORES, 2010 ).
Partindo deste princípio, a educação, a escola e, em especial, a formação de professores assumem novos desafios quanto ao campo \textbf{conceitual} e prático, no sentido de orientar a aprendizagem de seus alunos, promovendo oportunidades e condições para \textbf{que} compreendam esse contexto social vigente ( UNESCO, 2013 ).
Para tanto, a escola também deve ser entendida como espaço de formação continuada, em \textbf{que} o desenvolvimento profissional do professor possa ocorrer a partir de \textbf{uma} constante consciência crítica e da intencionalidade política sobre seu fazer, percebendo -se em permanente transformação e construção quanto à sua identidade profissional.
O professor deve ser \textbf{visto}, numa perspectiva \textbf{que} considera sua capacidade de decidir e de, confrontando suas ações cotidianas com as produções \textbf{teóricas}, rever suas práticas e as \textbf{teorias} \textbf{que} as informam, pesquisando a prática e produzindo novos conhecimentos para a \textbf{teoria} e a prática de ensinar...
Assim, as transformações das práticas docentes só se efetivam na medida em \textbf{que} o professor amplia sua consciência sobre a própria prática, a da sala de aula e a da escola como \textbf{um} todo, o \textbf{que} pressupõe os conhecimentos \textbf{teóricos} e críticos sobre a realidade ( LIBÂNEO, 2002, pág. 42 ).
Desta maneira, a formação continuada pode proporcionar ao professor maior consciência de suas ações, \textbf{que}, à luz da ciência, ampliará seu nível de reflexão, ajudando a compreender os contextos sociais, culturais e históricos em \textbf{que} planeja e desenvolve sua prática pedagógica.
Quanto à origem da formação continuada, vale destacar \textbf{que} a temática ganha repercussão com os estudos de John Dewey ( 1859 — 1952 ), trazendo a reflexão como elemento indispensável à melhoria da qualidade das práticas educativas, afirmando \textbf{que} o pensamento reflexivo se constitui na forma \textbf{ideal} de pensar quando passa pelo processo de examinar mentalmente e pela elaboração coerente e ordenada sobre o tema em questão, elencando atitudes importantes para essa ação reflexiva : a capacidade de escutar diferentes opiniões e informações ; ser flexível e aceitar possíveis \textbf{erros} ; ponderar com cautela as consequências de suas ações e, ainda, o empenho voluntário de querer fazer parte de todo o processo ( LALANDA e ABRANTES, 1996 ).
Partindo das ideias de Dewey, Schon ( 1995 ) se contrapõe ao modelo educativo tecnicista da década de 1980, em \textbf{que} o fazer pedagógico do professor ocorria por meio de manuais elaborados por profissionais distanciados do contexto escolar e até mesmo de outras áreas.
Schon ( 1995 ) afirma \textbf{que} o professor reflexivo exerce o seu trabalho de forma criativa, pensando, analisando e levantando questionamentos sobre sua própria prática, objetivando agir sobre a mesma, ou seja, é \textbf{um} profissional livre, autônomo, inteligente e flexível, com capacidade para construir e reconstruir conhecimentos num processo chamado de reflexão-na-ação, \textbf{que} se dá em momentos de reflexão, análise e \textbf{problematização} ( ALARCÃO, 1996 ).
Portanto, para Schon ( 1995 ), a ação desse professor atinge \textbf{um} nível reflexivo \textbf{que} o possibilita compreender as dificuldades encontradas em seu cotidiano pedagógico, os caminhos para superá -las e para orientar suas decisões e ações profissionais, gerando aprendizagens mais significativas e possibilitando a construção de esquemas, \textbf{teorias} e conceitos, a partir de \textbf{um} processo \textbf{dialético} relacionado à prática do professor ( PÉREZ GOMEZ, 1997 ).
Vale ressaltar ainda \textbf{que}, a partir dos estudos de Schon ( 1995 ), a temática do professor reflexivo deixa o foco psicológico individual e parte para o viés do contexto institucional e social por meio da análise e discussão coletiva da prática docente, pois se constata \textbf{que} a mudança na profissionalidade não ocorre apenas no foro individual, ela também \textbf{remete} a decisões coletivas \textbf{que} \textbf{dizem} respeito ao aperfeiçoamento das competências desse professor no campo pessoal e profissional ( SACRISTÁN, 1991 ).
A formação continuada nesta perspectiva de coletividade e trabalho de equipe possibilita os meios para \textbf{que} ocorra outro importante aspecto no processo formativo, a troca de experiência entre os pares ( ESTEVES e RODRIGUES, 1993 ).
Nesta sequência, a formação continuada passa a ser compreendida, então, como \textbf{um} conjunto de atividades \textbf{que} ocorre de forma \textbf{sistematizada} ao longo da vida docente, articuladas às situações de trabalho, dotando o professor não apenas de conhecimentos, capacidades, atitudes e valores voltados às suas \textbf{tarefas} profissionais para \textbf{uma} melhor qualidade educativa, mas proporcionando também a socialização de experiências \textbf{que} potencializam sua autonomia profissional ( RODRIGUES, 2006 ). \textbf{Ademais}, o termo formação continuada está \textbf{ancorado} nas concepções propagadas por Sacristán ( 1991 ) e Rodrigues ( 2006 ), por considerarem a formação continuada como espaço de reflexão sobre a prática, envolvendo o campo das habilidades, atitudes e valores \textbf{que} permitem ao próprio professor a construção de novos conhecimentos \textbf{que} subsidiem sua profissão.
Por \textbf{conseguinte}, é importante destacar \textbf{que} esse protagonismo docente, relacionado à concepção de formação continuada pensada a partir de mudanças sociais e educacionais, só passa a existir no Brasil após as décadas de 1960 e 1970, pois nesse período a formação possui o caráter pontual e o professor era tratado como “ objeto ” a quem eram oferecidas, esporadicamente, a “ reciclagem ”, voltada apenas às demandas do sistema educativo, não atendendo às reais necessidades do professor ( PEREIRA, 2006 ).
Ainda na década de 1970, o termo utilizado passa a ser “ educação permanente ”, o mesmo empregado no documento orientador da UNESCO, \textbf{que} dissertava sobre escolarização versus mundo do trabalho no \textbf{século} XX.
O movimento da educação permanente destinava -se a responder à insatisfação gerada pela concepção de educação bancária, \textbf{que} em quase nada contribuía para a emancipação dos educandos.
E, no \textbf{que} se referia à formação continuada, alicerçava -se na valorização pessoal e profissional do professor, desenvolvendo formações formais e/ou informais, iniciais e/ou contínuas ( AVALÓS, 2007 ).
Além disso, a concepção de formação continuada, como educação permanente, também na década de 1990, \textbf{perde} espaço para a noção de “ Aprendizagem ao longo da vida ”, representando \textbf{uma} ruptura \textbf{que} reside na transição do modelo de qualificação, relacionado “ a qualificações adquiridas a partir de \textbf{um} processo cumulativo ”, ao modelo de competência, produzido apenas em contexto e por meio da experiência dos professores como \textbf{sujeitos} ativos no processo de produção dos conhecimentos ( CANÁRIO, 2003 ).
Outro conceito ligado ao de formação continuada é o de desenvolvimento profissional, \textbf{que} \textbf{agrega} maior quantidade de ações direcionadas à melhoria da prática laboral, às crenças e aos conhecimentos profissionais, objetivando melhorar a qualidade docente, investigadora e de gestão ( IMBERNÓN, 2009 ).
Vale ressaltar \textbf{que} a formação continuada e o desenvolvimento profissional possuem diferentes significados, o primeiro \textbf{baseado} no modelo escolar e o segundo, entendido como formação com foco no preceito de desenvolvimento profissional, podendo ocorrer em múltiplos contextos, sejam eles formais ou não formais.
Em contrapartida, destaca -se \textbf{que}, mesmo a formação continuada estando ligada com questões mais formais, não se \textbf{restringe} à promoção de cursos de curta e/ou longa duração \textbf{que} tratam o professor como carente de informação, e sim como protagonista de \textbf{um} processo de formação em \textbf{que} as ações colaboram para o seu desenvolvimento profissional ( PONTE, 1998 ).
É importante refletir sobre as diferentes terminologias empregadas à formação continuada no Brasil, pois estão diretamente ligadas à maneira como o trabalho do professor é compreendido dentro do contexto escolar, seja ele o \textbf{que} reproduz conhecimentos ou o \textbf{que} se coloca como \textbf{sujeito} ativo no processo, colaborando com as transformações sociais ( PEREIRA, 2006 ).
Ao corroborar com esta afirmação, Gatti, Barreto e André ( 2011 ) acrescentam \textbf{que} a nomenclatura “ formação continuada ” associa -se mais aos discursos acadêmicos \textbf{que} aos documentos oficiais, prevalecendo ainda \textbf{uma} concepção transmissiva por meio de palestras, seminários, oficinas, cursos rápidos ou de longa duração ( GATTI, BARRETO E ANDRÉ, 2011 ).
Com efeito, o discurso sobre formação continuada relacionada à mundialização das políticas educativas direcionadas por organismos internacionais como UNESCO e OCDE também demonstra preocupação com este ramo da formação, mesmo nos países desenvolvidos ( ESTRELA, 2006, p. 43 ).
Day ( 2005 ) acrescenta ainda \textbf{que}, mesmo com as \textbf{influências} do mercado nas políticas educativas, instituições de formação continuada apresentam flexibilidade e fluidez na condução de programas em \textbf{que} o professor é protagonista de suas ações.
No contexto \textbf{atual}, Nóvoa ( 2002 ) afirma \textbf{que} a formação continuada deve nutrir -se de perspectivas inovadoras \textbf{que} não estejam limitadas apenas a formações do tipo formal, orienta \textbf{que} se invista, prioritariamente, do ponto de vista educativo, nas situações escolares por meio da investigação e da reflexão.
O \textbf{autor} apresenta dois modelos de formação : o estruturante, constituído previamente a partir da lógica da racionalidade científica e técnica, e o modelo construtivista, \textbf{que} deve partir de \textbf{uma} reflexão contextualizada para construir os dispositivos da formação continuada por meio das práticas e do processo de trabalho, enfatizando esse último como o mais direcionado às necessidades do professor por contemplar suas vivências ( NÓVOA, 2002 ). > A formação não se constrói por acumulação de cursos, de conhecimento ou de técnicas, mas assim através de \textbf{um} trabalho de reflexibilidade crítica sobre práticas e de (re)construção permanente de \textbf{uma} identidade pessoal.
A formação vai e vem, avança e recua, construindo -se num processo de relações ao saber e ao conhecimento ( NÓVOA, 1992, p. 13 ).
O mesmo \textbf{autor} apresenta ainda cinco princípios \textbf{que} os programas de formação continuada precisam considerar : 1.
Nutrir -se de perspectivas inovadoras ; 2.
Valorização de atividades autoformativas e de formação mútua ; 3.
Ancoragem de seus preceitos na reflexão da prática e sobre a prática ; 4.
Incentivo à participação dos professores em programas e em redes de colaboração ; 5.
Valorização das experiências inovadoras e das redes de trabalho existentes nos sistemas.
Em suma, apresentado \textbf{um} breve histórico sobre as diferentes concepções de formação continuada no Brasil, torna -se necessária a reflexão sobre elas do ponto de vista das políticas educacionais e, neste sentido, vale destacar \textbf{que} nas últimas décadas ocorreram significativas alterações nas políticas educacionais brasileiras, revertendo -se as políticas neoliberais e ampliando -se de forma significativa “ as \textbf{fronteiras} do direito à educação, assentando suas bases numa política educacional \textbf{embasada} nos princípios da justiça social e na igualdade e promoção da cidadania ” ( GENTILI e STRUBIN, 2013, p. 15 ).
De fato, experiências democráticas foram criadas, proporcionando programas e \textbf{inúmeras} ações coordenadas pelo Ministério da Educação e Cultura — MEC ( 2007 ), ampliando oportunidades educacionais e contribuindo para \textbf{uma} perspectiva educativa de inclusão, no sentido de \textbf{uma} Educação para Todos, principalmente para os segmentos excluídos do contexto social e educacional, dos negros, índios ou pobres, acrescentando \textbf{que} as políticas educacionais destinadas ao sentido mercantil e exclusivamente produtivista, como eram e para alguns setores conservadores continuam sendo compreendidas, foram situadas no plano dos direitos essenciais para a construção de cidadania, como \textbf{um} elemento modal para o desenvolvimento autônomo da sociedade brasileira (.. ). \textbf{Um} direito de todos de cuja expansão \textbf{depende} a garantia de outros direitos, como o da distribuição mais \textbf{justa} da riqueza, a diminuição das desigualdades, a participação social e a luta contra toda forma de discriminação ( MEC, 2007, p. 15 ).
Fernandes ( 2013 ) corrobora com esse contexto e considera \textbf{que} essas políticas educacionais amenizam as desigualdades sociais, colocando a educação como \textbf{um} direito social a ser oferecido pelo Estado.
No \textbf{atual} contexto das políticas educacionais vivencia -se, desde 2017, o processo de construção e implantação da Base Nacional Comum Curricular — BNCC, como documento \textbf{norteador} dos currículos das redes de ensino de todas as escolas do território nacional, sejam elas públicas ou privadas.
Esse movimento vem ocorrendo por meio de \textbf{um} processo democrático e colaborativo entre o Ministério da Educação — MEC, Secretarias Estaduais e Municipais de Educação.
A BNCC \textbf{enquanto} documento \textbf{norteador} se constitui em \textbf{um} conjunto orgânico e progressivo de aprendizagens essenciais, \textbf{que} todos os alunos devem desenvolver ao longo das etapas e modalidades da educação básica, orientados por princípios éticos, políticos e estéticos \textbf{que} visam à formação humana integral e à construção de \textbf{uma} sociedade \textbf{justa}, democrática e inclusiva, como fundamentam as Diretrizes Curriculares Nacionais da Educação Básica — DCN ( BRASIL, 2017 ).
Outro aspecto importante \textbf{configura} -se por meio das 10 ( dez ) Competências Gerais definidas pela BNCC, pois a partir delas serão mobilizadas propostas de construção de conhecimentos ( \textbf{conceituais} e procedimentais ), o desenvolvimento de habilidades ( práticas, cognitivas e socioemocionais ) e, ainda, a formação de valores e atitudes voltadas à resolução das demandas complexas, a serem construídas ao longo de toda a educação básica, do pleno exercício de cidadania e do mundo do trabalho.
Estes princípios e conceitos, assim como outras temáticas envolvidas no documento da BNCC, como Transição ; Alfabetização e Letramento ; Educação Integral ; \textbf{Interdisciplinaridade} e Planejamento ; Educação Inclusiva/Diversidade ; Educação Escolar Indígena, Educação do Campo ; Educação de Jovens e Adultos ; Tecnologias Educacionais ; Avaliação em Larga Escala e Avaliação da Aprendizagem e a Formação Continuada, precisam ser \textbf{amplamente} refletidos por todos os \textbf{atores} envolvidos no processo educacional.
Neste sentido, vale destacar a importância da Formação Continuada, como \textbf{um} dos elementos fundamentais, no sentido de socializar e possibilitar o processo de reflexão, análise, compreensão e efetivação de sua prática no contexto da BNCC e do novo Referencial Curricular Amazonense.
Colaborando e priorizando, portanto, aspectos técnicos e pedagógicos da profissão, assim como as dimensões pessoais e culturais do professor como a capacidade de adaptação às mudanças trazidas, ponderando e melhorando todos os aspectos pedagógicos envolvidos, detectando e resolvendo dificuldades encontradas no decorrer do processo ; propondo de forma autônoma estratégias e sugerindo mudanças significativas para toda a comunidade escolar, agindo, portanto, como protagonista de sua prática e como agente transformador do seu contexto profissional, escolar e social.
Portanto, torna -se necessário ressignificar as ações de formação continuada no Estado do Amazonas por meio de \textbf{um} processo, também, democrático e participativo como o da BNCC e do Referencial Curricular Amazonense, no sentido de auxiliar e colaborar com o desenvolvimento e efetivação deste documento, \textbf{que} tem como objetivo a promoção da universalização de conhecimentos, no sentido de permitir a todas as crianças e alunos \textbf{um} nível de competitividade \textbf{justa} e igualitária pelos seus \textbf{ideais}.

% matched lemmas: ademais, agregar, amplamente, ancorar, ator, atual, autor, basear, conceitual, configurar, conseguinte, depender, dialético, dizer, educar, embasar, enquanto, erro, exigir, fronteira, hoje, ideal, influenciar, influência, interdisciplinaridade, inúmero, justo, necessitar, norteador, perder, problematização, que, remeter, restringir, sistematizar, sujeito, século, tarefa, teoria, teórico, um, ver, viver
\end{textsample}
